% ==========================================
% CHAPITRE III : CONCEPTION ET DEVELOPPEMENT DU SYSTEME
% Titre et label declares dans main.tex (\label{chap:conception_developpement})
% Ton : recit de realisation au passe + present technique
% ==========================================

% Marqueur de travail \TODOchapo : defini dans main.tex.
% A supprimer de la version finale (rechercher "TODO" dans ce fichier).

\phantomsection
\section*{Introduction du chapitre}

Dans les chapitres précédents, nous avons présenté le cadre théorique et la démarche méthodologique. Ce chapitre décrit le système conçu pour donner accès aux connaissances scientifiques sur les plantes médicinales du Burkina Faso. Il présente OntoPlantMedBF, l'ontologie adaptée et peuplée, la base de connaissances orientée graphe, l'extraction des données, leur interrogation par le système RAG et l'application web.

% ============================================================
% III.1. ONTOLOGIE ADAPTÉE ET ENRICHIE
% ============================================================

\section{Ontologie adaptée et enrichie}
\label{sec:ontologie_adaptee}

En phase 1, nous avons sélectionné l'ontologie \textit{ontoMEDTRAD} \cite{kouame_systeme_2018} (section \ref{subsec:phase1} du chapitre \ref{chap:materiel_methodologie}) comme point de départ. Ce modèle constituait une base solide, mais incomplète pour le contexte burkinabè : il ne représentait pas les noms vernaculaires locaux, les niveaux de preuve ni la traçabilité détaillée des sources. En phase 4, nous en avons adapté le modèle conceptuel. En phase 5, nous avons enrichi l'ontologie avec les données et individus extraits de dix documents scientifiques. Enfin, en phase 6, nous avons importé dans Neo4j l'ontologie réajustée \textit{OntoPlantMedBF}.

\subsection{Vue d'ensemble du modèle}
\label{subsec:vue_ensemble_ontologie}

L'ontologie d'origine comptait 144 concepts, 76 propriétés associatives, 16 propriétés de données, 789 individus et 2 734 axiomes. Elle était formalisée en logique de description \textit{ALCHOIQ}, éditée sous Protégé~5 et vérifiée avec les raisonneurs HermiT et Pellet, qui n'y avaient signalé aucune incohérence. Nous en avons retenu une partie du schéma et l'avons étendue pour produire un modèle conceptuel adapté aux besoins de l'étude (phase 4). Ce modèle a ensuite été enrichi par des données et individus issus de dix documents scientifiques (phase 5), avant l'import et l'ajustement de cette représentation dans Neo4j (phase 6).

La figure \ref{fig:ontomedtrad_initiale} présente la structure de l'ontologie initiale, reconstituée à partir de la thèse.

\begin{figure}[H]
\centering
\begin{minipage}[t]{0.47\textwidth}
\centering
\includegraphics[width=\textwidth]{ontoMetrad-1.png}\\[1mm]
{\small (a) Partie supérieure de la hiérarchie}
\end{minipage}\hfill
\begin{minipage}[t]{0.47\textwidth}
\centering
\includegraphics[width=\textwidth]{ontoMetrad-2.png}\\[1mm]
{\small (b) Partie inférieure de la hiérarchie}
\end{minipage}
\caption{L'ontologie initiale ontoMEDTRAD reconstituée sous Protégé~5 (arbre des classes)}
\label{fig:ontomedtrad_initiale}
\end{figure}

Nous avons retenu d'ontoMEDTRAD les concepts relatifs aux plantes médicinales, à leurs parties, aux modes de préparation, aux maladies, à la localisation et aux références scientifiques. Nous les avons complétés par les classes et relations nécessaires au périmètre de notre étude, notamment les usages thérapeutiques, les noms vernaculaires, les niveaux de preuve, la traçabilité documentaire, la phytochimie et la toxicologie. Nous avons écarté les modules hors périmètre, dont la couche d'iconisation ontoICONE et les concepts relatifs aux praticiens, aux patients, à l'examen clinique et à l'anamnèse. L'iconisation étant un choix de représentation visuelle, son retrait ne remet pas en cause la réutilisation du modèle conceptuel. Le tableau \ref{tab:concepts_retenus_ecartes} détaille ces choix.

Nous avons réutilisé la structure conceptuelle, mais pas les 789 individus d'ontoMEDTRAD, associés au contexte et aux sources des travaux d'origine. En phase 4, nous avons défini le modèle conceptuel adapté et ses nouveaux éléments de schéma. En phase 5, nous avons enrichi l'ontologie avec les individus extraits de dix documents scientifiques et leurs liens aux sources. En phase 6, nous avons importé dans Neo4j l'ontologie réajustée OntoPlantMedBF.

\begin{table}[H]
\centering
\caption{Concepts retenus et écartés lors de l'adaptation}
\label{tab:concepts_retenus_ecartes}
\begin{tabular}{|L{4.6cm}|L{4.6cm}|L{3.8cm}|}
\hline
\textbf{Concepts retenus} & \textbf{Concepts écartés} & \textbf{Justification} \\
\hline
Plante, types et parties de plante & Iconisation (ontoICONE : Forme, Silhouette, Couleur\dots) & Représentation visuelle des plantes, hors du périmètre du système \\
\hline
Maladie (catégories CIM-10) & Praticiens (PMT, professions, expériences) & Gestion de la pratique des tradipraticiens, non utilisée \\
\hline
Mode de préparation & Patient, état du patient, anamnèse, examen clinique & Suivi clinique des patients, hors du périmètre \\
\hline
Références scientifiques & Indices (signes, symptômes) & Diagnostic clinique, non exploité \\
\hline
Localité, Pays & Glossaire, croyances, loi-cadre & Dimensions culturelles et réglementaires, non exploitées \\
\hline
Noms vernaculaires en langues locales & Langues (LangueLocales, LangueOfficielle) & Dimension linguistique reprise par les noms vernaculaires \\
\hline
\end{tabular}
\end{table}

Nous n'avons pas repris tels quels les concepts de recette, d'ingrédient et de remède d'ontoMEDTRAD. Nous avons retenu les classes \textit{UsageTherapeutique} et \textit{Preparation} pour représenter les besoins correspondants sous une forme adaptée. Nous n'avons pas conservé la classe générale \textit{Langues} comme axe autonome. Nous avons représenté la langue associée à un nom vernaculaire au moyen des propriétés de nom local et de code de langue. Ces concepts et relations structurent OntoPlantMedBF, dont la version peuplée contient les individus documentés dans les sources.

La figure \ref{fig:ontologie_conservee} présente le modèle réduit aux concepts retenus, importé dans Protégé~5. Cette étape intermédiaire a servi de base à la construction d'OntoPlantMedBF.

\begin{figure}[H]
\centering
\includegraphics[width=0.7\textwidth]{ontology-conserve.png}
\caption{Concepts d'ontoMEDTRAD conservés comme base d'OntoPlantMedBF}
\label{fig:ontologie_conservee}
\end{figure}

Le tableau \ref{tab:stats_ontologie} compare les métriques structurelles de l'ontologie d'origine à celles du modèle conceptuel adapté en phase 4 et de la version peuplée, dont le schéma a été complété lors de la phase 5 par des modes d'administration et de préparation, des propriétés de posologie et des parties de plante.

\begin{table}[H]
\centering
\caption{Statistiques de l'ontologie d'origine et du modèle conceptuel adapté}
\label{tab:stats_ontologie}
\begin{tabular}{|L{4.4cm}|C{2.6cm}|C{3.2cm}|C{3cm}|}
\hline
\textbf{Indicateur} & \textbf{Ontologie d'origine} & \textbf{Modèle adapté (phase 4)} & \textbf{Version peuplée (v4.0)} \\
\hline
Concepts (classes) & 144 & 69 (40 ajoutées) & 92 \\
\hline
Propriétés objet & 76 & 25 (20 ajoutées) & 29 \\
\hline
Propriétés de données & 16 & 20 (15 ajoutées) & 31 \\
\hline
Individus & 789 & Sans objet & 1~546 \\
\hline
Axiomes & 2 734 & Non relevé & \TODOchapo{relever dans Protégé : Active Ontology, Ontology metrics, Axiom} \\
\hline
\end{tabular}
\end{table}

Nous n'avons pas repris les 789 individus de l'ontologie d'origine dans le modèle conceptuel adapté. Le fichier \texttt{OntoPlantMedBF-peuplee.ttl} correspond à l'étape suivante : il enrichit ce modèle avec des individus extraits de dix documents scientifiques. Ces données ont ensuite été importées et ajustées dans Neo4j.

\begin{figure}[H]
\centering
\includegraphics[width=0.7\textwidth]{ontology-adaptee.png}
\caption{Modèle conceptuel adapté sous Protégé~5 (phase 4)}
\label{fig:protoge_ontologie}
\end{figure}

\subsection{Classes ajoutées}
\label{subsec:classes_ajoutees}

Nous avons construit OntoPlantMedBF en ajoutant au schéma initial des classes utiles à la représentation des connaissances ethnobotaniques et pharmacologiques étudiées. Ces ajouts complètent la sélection des concepts conservés. Le tableau \ref{tab:classes_ajoutees} en résume les principaux axes, sans lister toutes les classes.

\begin{table}[H]
\centering
\caption{Classes ajoutées à l'ontologie}
\label{tab:classes_ajoutees}
\begin{tabular}{|L{4.2cm}|L{5.4cm}|L{3.4cm}|}
\hline
\textbf{Ensemble de classes ajouté} & \textbf{Ce qu'il représente} & \textbf{Exemples de concepts} \\
\hline
Maladies et indications thérapeutiques & Catégories et indications du domaine & Concepts définis dans le schéma; données instanciées depuis les sources \\
\hline
Noms vernaculaires & Désignations locales et langue associée & Nom textuel et code de langue, si fournis \\
\hline
Localisation & Localités et régions pour structurer les données géographiques & Classes et individus issus des sources \\
\hline
Niveaux de preuve & Qualification telle que rapportée par la source & Traditionnel, \textit{in vitro}, \textit{in vivo}, clinique \\
\hline
Sources et traçabilité & Références et repères permettant de retrouver l'information & DOI, page, section et source \\
\hline
Préparations, parties et administration & Organes végétaux, procédés, formes et modes d'administration & Classes et individus selon les informations extraites \\
\hline
Phytochimie et sécurité & Composés, activités pharmacologiques et effets toxicologiques rapportés & Classes de composés et de risques; faits sourcés \\
\hline
\end{tabular}
\end{table}

Nous avons inséré les classes ajoutées dans les hiérarchies existantes lorsque cela était pertinent. Les parties végétales spécialisent \textit{PartieDePlante}, les formes galéniques et les procédés sont rattachés à \textit{Preparation} ou à \textit{ModeDePreparation}, et les catégories de métabolites spécialisent \textit{ComposeChimique}. Les classes \textit{NomVernaculaire}, \textit{NiveauDePreuve}, \textit{UsageTherapeutique} et \textit{ActivitePharmacologique} structurent les désignations locales, la qualification rapportée par les sources, les usages et les résultats pharmacologiques. Le modèle de phase 4 déclare 69 classes, dont 29 conservées et 40 ajoutées ; la version peuplée en déclare 92 et inclut des individus issus des sources.

\subsection{Propriétés ajoutées}
\label{subsec:proprietes_ajoutees}

Nous avons ajouté des propriétés de données pour décrire les noms vernaculaires, les usages, les niveaux de preuve, les activités pharmacologiques, les effets toxicologiques et les sources. La capture Protégé de la figure \ref{fig:proprietes_ajoutees} montre l'arborescence des propriétés de données. Elle comprend aussi des propriétés conservées de l'ontologie initiale ; le tableau \ref{tab:proprietes_ajoutees} présente des exemples d'ajouts.

\begin{table}[H]
\centering
\caption{Propriétés ajoutées à l'ontologie}
\label{tab:proprietes_ajoutees}
\begin{tabular}{|L{5.2cm}|L{7.8cm}|}
\hline
\textbf{Propriété} & \textbf{Rôle} \\
\hline
\textit{nomMooré}, \textit{nomDioula}, \textit{nomFulfuldé} & Noms vernaculaires de la plante dans les langues locales \\
\hline
\textit{estLocaliséeÀ} & Localisation géographique de la plante (vers \textit{Localite}) \\
\hline
\textit{nomVernaculaireTexte}, \textit{codeLangue} & Représentation structurée du nom local et de sa langue \\
\hline
\textit{niveauPreuve}, \textit{typeEtude} & Niveau de preuve rapporté par la source et type d'étude \\
\hline
\textit{source}, \textit{DOI}, page, section & Traçabilité documentaire de chaque connaissance \\
\hline
Dosage, toxicité & Précautions d'emploi documentées par les sources \\
\hline
\textit{descriptionActivite}, \textit{effetToxicologique} & Description des résultats rapportés \\
\hline
\end{tabular}
\end{table}

Les propriétés de données ajoutées sont des chaînes de caractères (\textit{xsd:string}). Nous avons créé \textit{nomMoore}, \textit{nomDioula} et \textit{nomFulfulde} pour décrire les noms locaux. Les propriétés \textit{nomVernaculaireTexte} et \textit{codeLangue} associent un nom vernaculaire à sa langue. \textit{niveauPreuve}, \textit{typeEtude}, \textit{source} et \textit{dosage} décrivent l'usage ; \textit{doi}, \textit{page} et \textit{section} identifient la référence. La propriété \textit{toxicite} consigne les précautions rapportées pour une plante. Les propriétés \textit{descriptionActivite} et \textit{effetToxicologique} décrivent les résultats rapportés. La localisation est décrite par \textit{estLocaliseeA}, du domaine \textit{Plante} à la classe \textit{Localite}, en complément de \textit{aCoordoneesGPS}, conservée de l'ontologie d'origine. Ces identifiants figurent dans le fichier Turtle.

\begin{figure}[H]
\centering
\includegraphics[width=0.96\textwidth]{proprietes_ajoutees.png}
\caption{Arborescence des propriétés de données dans Protégé}
\label{fig:proprietes_ajoutees}
\end{figure}

\subsection{Relations ajoutées}
\label{subsec:relations_ajoutees}

Nous avons ajouté des propriétés objet pour relier les plantes à leurs usages, les usages aux maladies cibles, parties utilisées, préparations et niveaux de preuve, ainsi que les plantes aux composés, activités et effets toxicologiques rapportés. Nous avons aussi structuré la provenance documentaire autour de la relation générale \textit{estDocumenteePar} et de ses sous-propriétés. La capture Protégé de la figure \ref{fig:relations_ajoutees} montre cette hiérarchie et comprend également des propriétés conservées. Le tableau \ref{tab:relations_ajoutees} présente les principales relations ajoutées.

\begin{table}[H]
\centering
\caption{Relations ajoutées à l'ontologie}
\label{tab:relations_ajoutees}
\begin{tabular}{|L{4cm}|L{9cm}|}
\hline
\textbf{Relation} & \textbf{Signification} \\
\hline
\textit{estDocumenteePar} & Relation générique rattachant une ressource ou affirmation à sa source scientifique \\
\hline
\textit{estDocumenteeDans} & Sous-propriété : relie une plante aux documents qui la décrivent \\
\hline
\textit{aPourNomVernaculaire} & Rattache une plante à ses désignations locales \\
\hline
\textit{estLocaliseeA} & Associe une plante à sa localisation géographique (\textit{Localite}) \\
\hline
\textit{provientDe}, \textit{aPourNiveauDePreuve} & Rattachent un usage à sa source et au niveau déclaré \\
\hline
\textit{aPourPartieUtilisee}, \textit{aPourPreparation} & Relient un usage à la partie végétale et à la préparation rapportées \\
\hline
\textit{aPourMaladieCible}, \textit{estRapportePour} & Relient l'usage à la maladie cible et la plante à l'usage documenté \\
\hline
\textit{contientCompose}, \textit{presenteActivite} & Représentent les composés et activités rapportés par les sources \\
\hline
\textit{presenteEffetToxicologique} & Relie une plante à un effet toxicologique rapporté \\
\hline
\end{tabular}
\end{table}

Nous avons associé le niveau de preuve à l'usage documenté, sans l'interpréter comme une validation indépendante de l'efficacité. La propriété textuelle \textit{niveauPreuve} conserve la valeur rapportée, et la relation \textit{aPourNiveauDePreuve} relie l'usage à la classe \textit{NiveauDePreuve}. La relation générale \textit{estDocumenteePar} rattache les ressources à une source scientifique. Nous l'avons spécialisée par \textit{provientDe} pour les usages, \textit{estDocumenteeDans} pour les plantes et des relations équivalentes pour les noms vernaculaires, les activités, les composés et les effets toxicologiques. Les individus n'emploient que deux niveaux de preuve, qui désignent des usages rapportés : nous les avons placés au rang 7 de l'échelle à sept niveaux, celui des savoirs traditionnels. Le modèle de phase 4 déclare 25 propriétés objet, dont 5 conservées et 20 ajoutées ; la version peuplée en déclare 29. Les relations \textit{aPourPartieUtilisee}, \textit{aPourPreparation} et \textit{aPourMaladieCible} décrivent le contenu d'un usage ; \textit{contientCompose}, \textit{presenteActivite} et \textit{presenteEffetToxicologique} décrivent les résultats rapportés. Ces liens structurent les informations, mais ne constituent pas une preuve d'efficacité clinique. La capture de la figure \ref{fig:relations_ajoutees} montre l'arborescence de provenance. Les domaines et portées sont déclarés dans le TTL et résumés au tableau \ref{tab:relations_ajoutees}.

\begin{figure}[H]
\centering
\includegraphics[width=0.96\textwidth]{relations_ajoutees.png}
\caption{Arborescence des propriétés objet dans Protégé}
\label{fig:relations_ajoutees}
\end{figure}

% ============================================================
% III.2. BASE DE CONNAISSANCES ORIENTÉE GRAPHE
% ============================================================

\section{Base de connaissances orientée graphe}
\label{sec:base_connaissances}

Nous avons utilisé Neo4j Desktop pour stocker et interroger les connaissances extraites selon le schéma de l'ontologie. Cette base représente les liens entre plantes, usages, propriétés et sources, que nous avons parcourus au moyen de requêtes. Nous décrivons le passage du modèle ontologique au graphe, son import et son interrogation.

\subsection{Du modèle ontologique au graphe}
\label{subsec:mapping_graphe}

Nous avons importé dans Neo4j l'ontologie réajustée OntoPlantMedBF. Chaque individu devient un nœud portant sa classe et ses superclasses, et chaque propriété objet devient une relation du même nom, conformément à \texttt{OntoPlantMedBF-peuplee.ttl}. Les propriétés de données deviennent des propriétés du nœud. Nous recopions aussi quelques valeurs, comme le nom scientifique ou l'identifiant de source, sous des noms plus simples pour l'affichage. Les individus et faits extraits sont intégrés avec leurs liens vers les sources. Le tableau \ref{tab:Mapping_noeuds} présente les principaux types de nœuds.

\begin{table}[H]
\centering
\caption{Correspondance entre les classes de l'ontologie et les nœuds du graphe}
\label{tab:Mapping_noeuds}
\begin{tabular}{|L{3.5cm}|L{4.3cm}|L{5.2cm}|}
\hline
\textbf{Nœud} & \textbf{Classe de l'ontologie} & \textbf{Propriétés principales} \\
\hline
\textit{Plante} & \textit{Plante} & Nom scientifique, famille, noms vernaculaires \\
\hline
\textit{Maladie} & \textit{Maladie} & Libellé repris de la source \\
\hline
\textit{UsageTherapeutique} & \textit{UsageTherapeutique} & Niveau de preuve rapporté, type d'étude, passage source, dosage \\
\hline
\textit{PartieDePlante} & \textit{PartieDePlante} et ses sous-classes & Partie utilisée \\
\hline
\textit{Preparation}, \textit{ModeDePreparation} & \textit{Preparation}, \textit{ModeDePreparation} et leurs sous-classes & Procédé, forme et ingrédients rapportés \\
\hline
\textit{SourceScientifique} & \textit{SourceScientifique} (classe ajoutée) & Titre, auteurs, année, DOI \\
\hline
\textit{ComposeChimique} & \textit{ComposeChimique} et ses sous-classes & Nom du composé \\
\hline
\textit{Localite} & \textit{Localite} et ses sous-classes, dont \textit{RegionBF} & Nom et coordonnées, si rapportées \\
\hline
\end{tabular}
\end{table}

Nous avons aligné les noms des relations du graphe sur les propriétés objet d'OntoPlantMedBF, notamment \textit{estRapportePour}, \textit{aPourMaladieCible}, \textit{aPourPartieUtilisee}, \textit{aPourPreparation}, les sous-propriétés de \textit{estDocumenteePar}, \textit{aPourNiveauDePreuve}, \textit{contientCompose}, \textit{presenteActivite}, \textit{presenteEffetToxicologique} et \textit{estLocaliseeA}. Nous avons repris les propriétés de données du TTL pour décrire les entités. Les noms Neo4j étant les identifiants OWL, la correspondance préserve le domaine et la portée définis dans l'ontologie, que l'import contrôle pour chaque relation.

\begin{table}[H]
\centering
\caption{Relations du graphe de connaissances}
\label{tab:relations_graphe}
\begin{tabular}{|L{4cm}|L{9cm}|}
\hline
\textbf{Relation} & \textbf{Sens retenu} \\
\hline
\textit{estRapportePour} & \textit{Plante} \(\rightarrow\) \textit{UsageTherapeutique} : rattache une plante à un usage documenté \\
\hline
\textit{aPourMaladieCible} & \textit{UsageTherapeutique} \(\rightarrow\) \textit{Maladie} : maladie ou indication rapportée \\
\hline
\textit{aPourPartieUtilisee} & \textit{UsageTherapeutique} \(\rightarrow\) \textit{PartieDePlante} : partie explicitement rapportée \\
\hline
\textit{aPourPreparation} & \textit{UsageTherapeutique} \(\rightarrow\) \textit{Preparation} : préparation ou procédé documenté \\
\hline
\textit{provientDe} & \textit{UsageTherapeutique} \(\rightarrow\) \textit{SourceScientifique} : source du fait rapporté \\
\hline
\textit{aPourNiveauDePreuve} & \textit{UsageTherapeutique} \(\rightarrow\) \textit{NiveauDePreuve} : niveau déclaré par la source \\
\hline
\textit{contientCompose}, \textit{presenteActivite}, \textit{presenteEffetToxicologique} & Liens vers les composés, activités et effets documentés \\
\hline
\end{tabular}
\end{table}

Ce mapping reprend les domaines et portées définis dans OntoPlantMedBF. Nous avons repris les noms OWL tels quels dans Neo4j afin de préserver cette équivalence lors du chargement et des requêtes.

\subsection{Import et peuplement}
\label{subsec:import_peuplement}

Nous avons validé les extractions JSON en conservant les formulations et les références des sources, puis les avons intégrées à l'ontologie sous forme d'individus. Nous avons ensuite écrit un programme d'import qui lit le fichier Turtle avec \texttt{rdflib}, contrôle les domaines et les portées, puis charge le graphe dans une transaction unique. Le chargement utilise \texttt{MERGE} sur l'URI de chaque individu : le relancer ne crée pas de doublon. Chaque fait reste relié à sa source et au passage d'origine. Les contradictions relevées dans les documents restent consignées dans les fichiers JSON, sans être tranchées. Nous n'avons ajouté au graphe aucune information absente des documents, à l'exception du rang attribué aux niveaux de preuve.

% Source : GraphRAG/application/ingestion/ttl_import.py (extrait)
\begin{quote}
\begin{verbatim}
# Création des nœuds, regroupés par ensemble d'étiquettes
tx.run(
    f"UNWIND $lignes AS ligne "
    f"MERGE (n:{LABEL_COMMUN} {{uri: ligne.uri}}) "
    f"SET n += ligne.props, n:{etiquettes}",
    lignes=lot,
).consume()

# Création des relations, regroupées par propriété objet
tx.run(
    f"UNWIND $lignes AS ligne "
    f"MATCH (a:{LABEL_COMMUN} {{uri: ligne.de}}) "
    f"MATCH (b:{LABEL_COMMUN} {{uri: ligne.vers}}) "
    f"MERGE (a)-[:{_sur(nom)}]->(b)",
    lignes=lot,
).consume()
\end{verbatim}
\end{quote}

Après l'import, le programme relit le graphe et compare ses comptes à ceux du fichier ; tout écart est signalé comme une erreur. Les résultats sont les suivants :

Le fichier compte 12~694 triplets. Le graphe contient 1~546 nœuds et 4~677 relations, comme attendu, sans anomalie de domaine ou de portée.

\begin{itemize}
    \item \textbf{Nœuds :} 291 préparations (dont 7 solvants ou excipients), 277 usages thérapeutiques, 198 composés chimiques, 192 noms vernaculaires, 160 maladies, 141 plantes, 109 genres, 61 familles, 46 activités pharmacologiques, 17 parties de plante, 12 localités, 11 modes d'administration, 11 modes de préparation, 10 sources, 6 effets toxicologiques, 2 ethnies et 2 niveaux de preuve.
    \item \textbf{Principales relations :} \textit{aPourPartieUtilisee} (417), \textit{aPourModeDePreparation} (414), \textit{aPourModeAdministration} (360), \textit{aPourPreparation} (284), \textit{estDocumenteePar} (284), \textit{estRapportePour}, \textit{provientDe} et \textit{aPourMaladieCible} (277 chacune), \textit{aPourNiveauDePreuve} (270), \textit{estDocumenteeDans} (230), \textit{estLocaliseeA} (204).
\end{itemize}

\begin{figure}[H]
\centering
\fbox{\parbox[c][4.5cm][c]{0.8\textwidth}{\centering
\TODOchapo{Capture Neo4j Browser d'un nœud Plante et de ses voisins, à enregistrer sous \texttt{images/capture\_neo4j.png}.}\\[2mm]
Vue du graphe peuplé dans Neo4j Browser (extrait du sous-graphe d'une plante)}
}
\caption{Le graphe de connaissances peuplé}
\label{fig:graphe_neo4j}
\end{figure}
% Une fois la capture faite, remplacer tout le cadre \fbox{...} ci-dessus par :
% \includegraphics[width=0.8\textwidth]{capture_neo4j.png}

\subsection{Interrogation de la base}
\label{subsec:interrogation_base}

Nous avons défini quatre familles de requêtes en phase 6 (section \ref{subsec:phase6} du chapitre \ref{chap:materiel_methodologie}) : requêtes simples, requêtes complexes, requêtes de traçabilité et requêtes spatiales. Nous en donnons des exemples ci-dessous.

\begin{quote}
\begin{verbatim}
// 1. Requête simple : plantes documentées contre le paludisme
MATCH (p:Plante)-[:estRapportePour]->(u:UsageTherapeutique)
      -[:aPourMaladieCible]->(m:Maladie)
WHERE toLower(m.nom) CONTAINS 'paludisme'
   OR toLower(m.nom) CONTAINS 'malaria'
RETURN DISTINCT p.nomScientifique AS plante, u.usageId AS usage

// 2. Requête complexe : recettes associant plusieurs plantes
MATCH (p:Plante)-[:estRapportePour]->(u:UsageTherapeutique)
      -[:aPourPreparation]->(:Preparation)
      -[:associePlante]->(a:Plante)
RETURN DISTINCT p.nomScientifique AS plante,
       a.nomScientifique AS associee, u.usageId AS usage

// 3. Requête de traçabilité : sources d'une plante
MATCH (:Plante {nomScientifique: 'Annona senegalensis'})
      -[:estDocumenteeDans]->(s:SourceScientifique)
RETURN s.sourceId AS source

// 4. Requête spatiale : plantes d'une localité
MATCH (p:Plante)-[:estLocaliseeA]->(:Localite {nom: 'Pays San'})
RETURN DISTINCT p.nomScientifique AS plante
\end{verbatim}
\end{quote}

Ces requêtes, exécutées sur la base peuplée, renvoient respectivement 16, 43, 6 et 93 lignes (tableau \ref{tab:volumes_requetes} du chapitre \ref{chap:resultats_discussion}).

Pour répondre à la question « quelles plantes sont rapportées contre le paludisme au Pays San ? », une requête relie les nœuds \textit{Plante}, \textit{UsageTherapeutique}, \textit{Maladie} et \textit{Localite}, puis renvoie les noms scientifiques correspondants, ici 8 plantes. Nous avons vérifié le fonctionnement de la base et les résultats des requêtes avant d'intégrer le système d'accès en langage naturel.

% ============================================================
% III.3. INTÉGRATION DE L'INTELLIGENCE ARTIFICIELLE
% ============================================================

\section{Intégration de l'intelligence artificielle}
\label{sec:integration_ia}

Nous avons utilisé des modèles de langue à deux étapes. Nous les avons d'abord sollicités pour extraire et structurer les connaissances du corpus, puis mis en place un système RAG pour interroger le graphe en langage naturel. Cette section décrit l'extraction, le système d'accès et la traçabilité entre les sources et les réponses.

\subsection{Extraction assistée des connaissances}
\label{subsec:extraction_llm}

\subsubsection{Environnement d'extraction}

Nous avons réalisé l'extraction avec le modèle \textbf{gpt-6-luna} d'OpenAI, accessible par l'environnement OpenWork, en soumettant les documents un par un. L'extraction a porté sur dix des 35 documents du corpus, les 19 et 20 octobre 2026. Le document de Ky et al., numérisé sans couche de texte, a dû passer par une reconnaissance optique de caractères, et son extraction reste partielle.

\subsubsection{Le prompt d'extraction}

Nous avons adressé au modèle une même consigne pour chaque document. Longue d'environ 10~700 caractères, elle précise les règles à respecter, les éléments à extraire et le format de sortie. En voici le début :

\begin{quote}
\begin{verbatim}
Tu es un assistant d'extraction de connaissances scientifiques
spécialisé dans l'ethnobotanique, l'ethnopharmacologie et la
phytochimie des plantes médicinales du Burkina Faso.

OBJECTIF
À partir du document fourni, extrais les informations explicitement
présentes et organise-les selon le schéma de connaissances décrit
ci-dessous, compatible avec notre ontologie adaptée ontoMEDTRAD.

RÈGLES IMPÉRATIVES
1. Utilise uniquement les informations présentes dans le document.
2. N'invente, ne déduis et ne complète aucune valeur manquante.
3. Distingue strictement :
   - un usage traditionnel rapporté ;
   - une activité pharmacologique observée en laboratoire ;
   - un résultat sur un modèle animal ;
   - un résultat clinique chez l'humain.
[...]
\end{verbatim}
\end{quote}

Un champ absent du document reçoit la valeur \texttt{null}, ou une liste vide. La consigne interdit de compléter une valeur par déduction et demande de rattacher chaque fait à sa page, sa section ou son tableau, et de signaler les contradictions sans les trancher. Ces règles limitent l'ajout d'éléments non étayés dans les fiches.

\subsubsection{Format de sortie}

La sortie est un document JSON comprenant la description du document, les plantes, les contradictions relevées, les informations absentes et un contrôle de l'extraction. Voici un extrait réel, tiré de la fiche d'\textit{Acacia albida} dans l'extraction de Ky et al. :

\begin{quote}
\begin{verbatim}
{
  "nom_scientifique_source": "Acacia albida Del.",
  "famille_botanique": "Mimosaceae",
  "synonymes": [],
  "noms_vernaculaires": [],
  "localisations": [{
    "localite": "Zone d'étude comprenant les communes de
                 Bissiga, Lalgaye et Tenkodogo",
    "province": null,
    "region": "Centre-Est",
    "coordonnees": null,
    "reference_passage": {
      "page": "1292",
      "section": "Appendix 1",
      "citation_courte": "Acacia albida Del."
    }
  }],
  "usages_therapeutiques": [{
    "indication_source": "Articular pains",
    "parties_utilisees": ["écorce"],
    ...
  }]
}
\end{verbatim}
\end{quote}

\subsubsection{Déroulement et contrôle}

Nous avons traité les documents un par un. Une première version de la consigne, jugée trop longue, a été raccourcie avant l'extraction. Un contrôle automatique en lecture seule a ensuite vérifié les fichiers : les 10 fichiers JSON sont valides et contiennent 230 fiches de plantes, 277 usages, 46 activités, 51 entrées de composés et 6 observations toxicologiques, sans aucun fait dépourvu de référence au passage source. Nous n'avons pas chronométré l'extraction.

\subsection{Système RAG d'accès aux connaissances}
\label{subsec:systeme_rag}

Nous avons mis en place un système de génération augmentée par récupération (RAG, \textit{Retrieval-Augmented Generation}). Il associe la recherche dans le graphe à la génération de réponses par un modèle de langue \cite{lewis2020rag}. Son architecture relie l'utilisateur, Streamlit, Neo4j et Llama 3 (figure \ref{fig:architecture_rag}). Nous avons confié la sélection des données à des requêtes fixes plutôt qu'au modèle : une requête générée par un modèle peut se tromper sans que l'utilisateur le voie, alors que des requêtes fixes garantissent que les faits proviennent bien du graphe. Llama 3 ne fait que rédiger la réponse à partir de ces faits.

\begin{figure}[H]
\centering
\includegraphics[width=\textwidth]{architecture_rag.png}
\caption{Chaîne de traitement du système RAG}
\label{fig:architecture_rag}
\end{figure}

\subsubsection{Les six étapes de la chaîne de traitement}

Nous avons organisé le traitement des requêtes en six étapes :

\begin{enumerate}[label=\textbf{\arabic*.}, leftmargin=1.2cm]
    \item \textbf{Réception de la question.} L'utilisateur formule une question en langage naturel, par exemple « Quelles plantes sont utilisées contre le paludisme ? ». L'application la nettoie et la transmet à la couche RAG.
    \item \textbf{Détection des entités.} Le système compare la question au vocabulaire réel du graphe : plantes, noms vernaculaires, parties, modes de préparation, composés et localités. Les libellés de maladies reprenant la langue de chaque source, une table d'équivalents relie par exemple « paludisme » à « malaria ».
    \item \textbf{Interrogation de Neo4j.} Chaque entité détectée déclenche une requête Cypher fixe et paramétrée, en lecture seule. Si aucune entité n'est reconnue, le modèle propose en repli une requête Cypher, que le système valide en lecture seule avant de l'exécuter.
    \item \textbf{Assemblage du contexte.} Les résultats sont réunis en faits structurés, chacun conservant son niveau de preuve et ses sources.
    \item \textbf{Génération de la réponse.} Llama 3 reçoit la question, les faits retrouvés et l'historique récent, et rédige une réponse qui ne doit s'appuyer que sur eux. Si aucun fait n'est retrouvé, le système indique que l'information n'est pas disponible, sans demander de réponse au modèle.
    \item \textbf{Affichage des sources.} L'interface présente la réponse, puis les références scientifiques correspondantes (auteur, année, DOI) et le niveau de preuve de chaque usage.
\end{enumerate}

\subsubsection{Câblage des composants}

Nous avons utilisé les composants LangChain présentés au tableau \ref{tab:langchain} du chapitre \ref{chap:materiel_methodologie}. \texttt{PromptTemplate} porte les consignes et \texttt{ChatOpenAI} appelle Llama 3 par l'API d'Ollama, exposée par ngrok, comme indiqué au tableau \ref{tab:llama_acces}. Le pilote \texttt{neo4j} assure l'accès à la base. L'adresse du modèle est lue dans la configuration, et le fichier qui la contient n'est pas versionné. L'extrait suivant montre l'instanciation du modèle.

% Source : GraphRAG/application/models/llm.py (extrait)
\begin{quote}
\begin{verbatim}
if config.utilise_ollama:
    from langchain_openai import ChatOpenAI

    # Ollama expose une API compatible OpenAI sur /v1 ;
    # sa clé est ignorée.
    return ChatOpenAI(
        model=modele,  # llama3:8b
        base_url=f"{config.ollama_url}/v1",
        api_key="ollama",
        temperature=config.temperature,  # 0
        timeout=180,
        max_retries=1,
    )
\end{verbatim}
\end{quote}

Le notebook Colab installe Ollama, télécharge le modèle, démarre le serveur avec un contexte de 8~192 jetons, puis ouvre le tunnel ngrok. Nous avons vérifié pendant l'évaluation que le modèle était entièrement chargé sur le GPU.

\subsection{Traçabilité des réponses}
\label{subsec:tracabilite_reponses}

Nous avons maintenu la traçabilité à chaque étape. L'extraction ne conservait que les informations rapportées par les sources et laissait vides les champs absents. Lors de l'import, chaque connaissance est restée reliée à un nœud \textit{SourceScientifique}. Les requêtes Cypher permettaient de retrouver ces références, affichées avec les réponses. Nous avons repris les niveaux de preuve tels quels, sans évaluation indépendante, et aucun expert externe n'a relu les extractions. Ces limites sont discutées au chapitre \ref{chap:resultats_discussion}.

% ============================================================
% III.4. DÉVELOPPEMENT DE L'APPLICATION WEB
% ============================================================

\section{Développement de l'application web}
\label{sec:application_web}

Nous avons développé une application web avec Streamlit pour donner accès au système. Elle recueille les questions, appelle la chaîne RAG et présente les réponses avec leurs sources. Les raisons du choix de Streamlit sont exposées au chapitre \ref{chap:materiel_methodologie}, section \ref{subsec:environnement}.

\subsection{Architecture de l'application}
\label{subsec:architecture_app}

Nous avons organisé l'application autour d'un point d'entrée, de huit écrans et d'un état de session partagé : accueil, assistant, recherche, plantes, pathologies, graphe, sources et à propos. Chaque question posée à l'assistant appelle la couche RAG, qui renvoie une réponse et les sources à afficher.

\subsection{Fonctionnalités}
\label{subsec:fonctionnalites}

Le tableau \ref{tab:fonctionnalites} récapitule les fonctions définies en phase 7 (section \ref{subsec:phase7}).

\begin{table}[H]
\centering
\caption{Fonctionnalités de l'application web}
\label{tab:fonctionnalites}
\begin{tabular}{|L{3.6cm}|L{7.4cm}|L{2cm}|}
\hline
\textbf{Fonctionnalité} & \textbf{Description} & \textbf{État} \\
\hline
Chatbot & Question en langage naturel, réponse générée par le système RAG, sources affichées & Terminé \\
\hline
Recherche par plante & Consultation d'une fiche à partir du nom scientifique ou vernaculaire & Terminé \\
\hline
Recherche par pathologie & Liste des plantes documentées contre une maladie donnée & Terminé \\
\hline
Fiches descriptives & Usages, parties utilisées, modes de préparation, niveaux de preuve et références & Terminé \\
\hline
Affichage des sources & Références complètes (auteur, année, DOI) sous chaque information & Partiel \\
\hline
Visualisation du graphe & Représentation des relations entourant une plante ou une maladie & Partiel \\
\hline
\end{tabular}
\end{table}

Deux fonctions restent partielles. Les références s'affichent sous la réponse et non sous chaque information, et la visualisation du graphe part d'une plante, pas d'une maladie.

\subsection{Gestion de session et déploiement}
\label{subsec:session_depl}

Nous avons conservé l'historique des échanges dans l'état de session de Streamlit et l'avons transmis au modèle dans la consigne de génération. L'historique sert à la rédaction, mais pas à la détection des entités : une question de suivi qui ne nomme pas son sujet ne déclenche donc pas de requête fixe. L'application fonctionne sur un poste local sous Windows~11 et accède au modèle par l'API exposée avec ngrok ; seule l'API du modèle est exposée, pas l'application.

\begin{quote}
\begin{verbatim}
python -m streamlit run application\app.py
\end{verbatim}
\end{quote}

L'application s'ouvre sur le port 8501. La session Colab n'est active que pendant son utilisation et son adresse change à chaque lancement. Il n'existe pas de mise en ligne pérenne.

\begin{figure}[H]
\centering
\fbox{\parbox[c][4.5cm][c]{0.8\textwidth}{\centering
\TODOchapo{Deux captures à enregistrer sous \texttt{images/ecran\_assistant.png} (une réponse avec ses sources) et \texttt{images/ecran\_fiche.png} (fiche d'une plante).}\\[2mm]
Assistant avec une réponse sourcée et fiche descriptive d'une plante}
}
\caption{Écrans principaux de l'application web}
\label{fig:ecrans_application}
\end{figure}
% Une fois les captures faites, remplacer tout le cadre \fbox{...} ci-dessus par :
% \begin{minipage}[t]{0.48\textwidth}\centering
% \includegraphics[width=\textwidth]{ecran_assistant.png}\\ Assistant
% \end{minipage}\hfill
% \begin{minipage}[t]{0.48\textwidth}\centering
% \includegraphics[width=\textwidth]{ecran_fiche.png}\\ Fiche avec sources
% \end{minipage}

% ============================================================
% CONCLUSION DU CHAPITRE
% ============================================================

\section*{Conclusion du chapitre}
\addcontentsline{toc}{section}{Conclusion du chapitre}

Ce chapitre a présenté le système issu de la démarche du chapitre \ref{chap:materiel_methodologie}. Nous avons construit OntoPlantMedBF à partir d'ontoMEDTRAD, l'avons adaptée au contexte burkinabè et peuplée avec des connaissances sourcées ; son schéma organise également la base Neo4j. Nous avons décrit l'extraction assistée par modèle de langue, le système RAG qui interroge le graphe à partir des questions et l'application web qui présente les réponses. Nous avons rattaché les connaissances à leurs sources, laissé vides les champs absents et repris les niveaux de preuve rapportés par les documents. Au chapitre \ref{chap:resultats_discussion}, nous avons évalué les extractions, les réponses et l'interface, puis discuté les résultats et leurs limites.
